\documentclass[10pt,a4paper]{amsart}
\usepackage[margin=19mm]{geometry}
\usepackage{amsmath,amssymb,mathtools}
\usepackage[scr=rsfs,cal=euler]{mathalfa}
\usepackage{xfrac}
\usepackage[unicode,hidelinks,bookmarks=false]{hyperref}
\newcommand{\ie}{i.e.\ }
\newcommand{\cf}{cf.\ }
\newcommand{\resp}{respectively\ }
\newcommand{\Subsection}{\subsection}
\providecommand{\nobreakdash}{\nobreak\hbox{-}}
\setlength{\emergencystretch}{2em}
\multlinegap0pt
\usepackage{amssymb}
\newtheorem{lem}{Lemma}
\newcommand{\be}{\begin{equation}}
\newcommand{\bea}{\begin{eqnarray}}
\renewcommand{\d}{\displaystyle}
	\newcommand{\ee}{\end{equation}}
	\newcommand{\eea}{\end{eqnarray}}
\renewcommand{\leq}{\leqslant}\renewcommand{\geq}{\geqslant}
\newcommand{\ol}{\overline}
\title{A GENERAL CHARACTERIZATION ON THE UNIQUENESS PROBLEM OF L-FUNCTIONS AND GENERAL MEROMORPHIC FUNCTIONS}
\author{Sanjay Mallick, Ripan Saha}
\date{Source: arXiv:2512.17205v1 (CC BY 4.0)}
\begin{document}
\maketitle

\noindent\emph{Source and licence.} A GENERAL CHARACTERIZATION ON THE UNIQUENESS PROBLEM OF L-FUNCTIONS AND GENERAL MEROMORPHIC FUNCTIONS. Version arXiv:2512.17205v1 (CC BY 4.0), distributed by arXiv under the Creative Commons Attribution 4.0 International licence, http://creativecommons.org/licenses/by/4.0/, as recorded in the arXiv licence metadata for this exact version and shown on the abstract page https://arxiv.org/abs/2512.17205v1. Full licence notice: \texttt{licenses/arxiv-2512.17205-CC-BY-4.0.txt}.\par\smallskip

\noindent\textit{Editorial scope.} Counted unit: the article's lem lem (source0.tex lines 570--576). The article's complete original proof of it is delivered with it (source0.tex lines 577--592, one \texttt{proof} environment of 16 lines). No mathematics is written, repaired or re-derived: the statement and the proof are the article's own lines, in the article's own notation, and no retained line is substituted or re-flowed.  No further source-local context is needed: every cross-reference of every retained line resolves inside the two delivered spans.  Nothing was omitted from any retained span, so the package carries no omission record.  No retained line contains a citation, so the wrapper carries no bibliography and no bibliography entry.  No retained line contains a figure environment, an image-inclusion command or drawing code, so no image is delivered and the package declares no asset dependency.  Editorial formatting only, in addition: an AMS \texttt{amsart} preamble that restates the article's own declarations and macros used by the retained lines, namely the environments \texttt{lem} on separate counters, together with the standard packages those lines need.  The retained environments print in this document as Lemma 1 in order of appearance, whereas the article prints the same items with the numbering of its own full document.  The complete native source of the article as distributed in the arXiv e-print for this exact version is bundled unchanged as \texttt{sources/191301/source0.tex}.
\begin{lem}
Let $F^{*}-1=\d\frac{a_{n}\prod\limits_{i=1}^{n}(f-w_{i})}{\psi(f)}$ and $G^{*}-1=\d\frac{a_{n}\prod\limits_{i=1}^{n}(\mathcal{L}-w_{i})}{\psi(\mathcal
{L})}$, where $f$  be a non-constant meromorphic function, $\mathcal{L}$ be an non-constant L-function, $a_{n},w_{i}\in\mathbb{C}-\{0\}$; $\forall i\in\{1,2,\ldots,n\}$ and  $\psi(z)$ be a polynomial of degree less than $n$ with $\psi(w_{i})\neq0$; $\forall i\in\{1,2,\ldots,n\}$. Further suppose that $F^{*}$ and $G^{*}$  share $(1,t)$, where $t\in\mathbb{N}\cup\{0\}$.  Then 
\be\nonumber \ol{N}_L(r,1;F^{*})\leq \frac{1}{t+1}\left[\ol{ N}(r,\infty;f)+\ol{ N}(r,0;f)-N_1(r,0,f')\right]+O(\log r),\ee
where $N_1(r,0,f')=N(r,0;f'|f\ne 0,w_1,w_2,...,w_n).$
Similar expression also holds for $\ol  N_L(r,1;G^{*})$.
\end{lem}

\begin{proof}
Since $F^{*}$ and $G^{*}$  share $(1,t)$, so using the first fundamental theorem we find that
\bea  \ol{ N}_L(r,1;F^{*}) &\leq& \ol{ N}(r,1;F^{*}|\geq t+2)\nonumber\\&&
\leq \frac{1}{t+1}\left[ N(r,1;F^{*})-\ol{ N}(r,1;F^{*})\right]
\nonumber\\&&
\leq \frac{1}{t+1}\left[\sum_{i=1}^{n}\left(N(r,w_i;f)-\ol{ N}(r,w_i;f)\right)\right]	\nonumber\\&&
\leq \frac{1}{t+1}\left[N(r,0;f'|f\ne 0)-N_1(r,0,f')\right]\nonumber\\&&
\leq \frac{1}{t+1}\left[N(r,0;\frac{f'}{f})-N_1(r,0,f')\right]\nonumber\\&&
\leq \frac{1}{t+1}\left[N(r,\infty;\frac{f}{f^{'}})-N_1(r,0,f')\right]+O(\log r)\nonumber\\&&
\leq \frac{1}{t+1}\left[N(r,\infty;\frac{f^{'}}{f})-N_1(r,0,f')\right]+O(\log r)\nonumber\\&&
\leq \frac{1}{t+1}\left[\ol{ N}(r,\infty;f)+\ol{ N}(r,0;f)-N_1(r,0,f')\right]+O(\log r)\nonumber\
%\\&&
%	\leq \frac{1}{p+1}\left[\ol{ N}(r,0;f)-N_1(r,0,f')\right]+S(r,f).\;\;[\because \ol{ N}(r,\infty;f)=O(\log r)]
\eea %as $f$ has finitely many poles and hence $\ol N(r,\infty;f)=S(r,f)$.
This proves the lemma.
\end{proof}

\end{document}
